\begin{definition}
  Let $X$ be a type. A \textbf{relation} on $X$ is a map $X \to X \to \Prop$. 
  Such a relation is called \textbf{reflexive} whenever it contains the diagonal $\Delta$ as a subset.
  For $R$ a relation, we denote $R^{-1}$ for the relation $\lambda x y . R(y,x)$. If $R = R^{-1}$, we call $R$ \textbf{symmetric}.
  For $R,S$ relations, we denote $R \circ S$ for the relation $\lambda x z . \exists_{y:X} R(x,y) \wedge S(y,z)$. 
\end{definition}
\begin{remark}
  We will silently be using that if $R$ is a symmetric relation and $R(x,y)$ and $R(z,y)$ then $(R\circ R)(x,z)$. 
  We will also use silently that whenver $R$ is a reflexive relation, $R\cap R^{-1}$ is a symmetric reflexive relation contained in $R$. 
\end{remark}
